\section{Clustered Spectra and Solver Cost}
\label{sec:spectra}

The chord law in \cref{thm:sublevel-chord-law} says when the condition
number must diverge, but not how the eigenvalues between its two edges are
distributed.  This section resolves that question for the logarithmic
barrier of a degenerate LP and then asks what the resulting spectrum costs.
The answer separates solver families.  The spectrum eventually consists of
two intervals with parameter-independent internal condition numbers, so a
classical Krylov method needs only polylogarithmically many matrix--vector
products.  Quantum polynomial solvers cannot use the empty spectral gap in
the same way: in sufficiently large dimension, plain block-encoding access
costs $\Thetatilde(\kappa)$ in the worst case, whereas the factor access already
defined in \cref{def:plain-factor-access} has a
$\Thetatilde(\sqrt\kappa)$ polynomial route under norm-tight encodings.
Clustering improves neither quantum model beyond the dependence its access
model already provides, while it exposes a classical--quantum
separation on the Newton systems themselves.  There is an important
qualification.  At exact central points the path's own Newton right-hand
sides are asymptotically orthogonal to the weak modes, so path-following
steps can use a condition-number-free window solve after the gap opens.

\subsection{The two-cluster tail}
\label{sec:spectra-two-cluster}

Consider the standard-form LP \eqref{eq:lp-pair} and suppose that its primal
feasible polytope $P$ is compact with $P^\circ\ne\varnothing$.  Let
$Z\in\R^{n\times p}$ have orthonormal columns spanning $\ker A$, assume as
in \cref{sec:conditioning} that the objective is nonconstant on the affine
hull, and follow
the exact log-barrier central path $(x(\mu),y(\mu),s(\mu))$.  Let $B$ be the
maximal primal optimal support and $N=[n]\setminus B$; equivalently, this is
the Goldman--Tucker optimal partition.  Write $P_B$ and $P_N$ for coordinate
projections, and put $D_P=\operatorname{diam}P$.

Choose a fixed central-path tail $0<\mu\leq\bar\mu$ on which
\begin{equation}
 S_*=\sup_{0<\mu\leq\bar\mu}\max_j s_j(\mu)<\infty,
 \qquad
 \beta=\inf_{0<\mu\leq\bar\mu}\min_{j\in B}x_j(\mu)>0.
 \label{eq:spectra-tail-constants}
\end{equation}
Strict complementarity of the central-path limit also gives
$\underline s_N:=\inf_{0<\mu\leq\bar\mu}\min_{i\in N}s_i(\mu)>0$.
These are instance constants, not asymptotic variables.

\begin{proposition}[Two-cluster log-barrier Hessian]
\label{prop:two-cluster-hessian}
Let
\begin{equation}
 H_\mu=Z^\top X(\mu)^{-2}Z=H_{\rm red}(\mu),
 \qquad r=\rank(P_NZ),
 \label{eq:two-cluster-hessian}
\end{equation}
where the equality uses \cref{sec:newton-formulations}'s reduced-Hessian notation
with $W=Z$.  If $r>0$, let $\sigma>0$ be the smallest positive singular
value of $P_NZ$.  The largest $r$ eigenvalues of $H_\mu$ then lie in
\begin{equation}
 \left[
  \frac{\sigma^2\min_{i\in N}s_i(\mu)^2}{\mu^2},
  \frac{\max_{i\in N}s_i(\mu)^2}{\mu^2}+\beta^{-2}
 \right].
 \label{eq:two-cluster-top}
\end{equation}
If $p-r>0$, the remaining $p-r$ eigenvalues lie in
\begin{equation}
 [D_P^{-2},\,\beta^{-2}].
 \label{eq:two-cluster-bottom}
\end{equation}
Each nonempty interval has a $\mu$-independent internal ratio.  When both
are nonempty, their separation grows as $\Theta(\mu^{-2})$; when $r=0$ the
top-cluster assertion is empty.
\end{proposition}

\begin{proof}
Complementarity gives $x_i(\mu)^{-2}=s_i(\mu)^2/\mu^2$.  Split
\begin{equation}
 H_\mu=H_{\mu,N}+H_{\mu,B},
 \quad
 H_{\mu,N}=Z^\top P_N^\top X_N^{-2}P_NZ,
 \quad
 H_{\mu,B}=Z^\top P_B^\top X_B^{-2}P_BZ.
 \label{eq:two-cluster-split}
\end{equation}
The first summand is positive semidefinite of rank $r$.  When $r>0$, its nonzero
eigenvalues are the squared nonzero singular values of $X_N^{-1}P_NZ$ and
therefore lie in
\[
 \left[
  \frac{\sigma^2\min_{i\in N}s_i(\mu)^2}{\mu^2},
  \frac{\max_{i\in N}s_i(\mu)^2}{\mu^2}
 \right].
\]
The tail bound on $x_B$ gives
$0\preceq H_{\mu,B}\preceq\beta^{-2}I$.
Weyl monotonicity now gives the lower edge in
\eqref{eq:two-cluster-top}, and Weyl's upper inequality gives its upper
edge.  Since $\ker H_{\mu,N}=\ker(P_NZ)$ has dimension $p-r$, the min--max
principle applied to this subspace gives
$\lambda_{p-r}^{\uparrow}(H_\mu)\leq\beta^{-2}$ when $p-r>0$, where
eigenvalues are ordered increasingly.  Hence every one of the lowest $p-r$
eigenvalues has the claimed upper bound.

It remains to bound their lower edge.  For every Euclidean unit tangent
direction, one of its two signs exits the compact polytope within distance
$D_P$.  The Dikin containment argument in
\cref{thm:lambda-min-rigidity} therefore gives
$H_\mu\succeq D_P^{-2}I$.  Finally,
\eqref{eq:spectra-tail-constants} and $\underline s_N>0$ bound each
nonempty interval's endpoint ratio independently of $\mu$.  If both clusters
are nonempty, the lower endpoint of \eqref{eq:two-cluster-top} is
$\Theta(\mu^{-2})$ while \eqref{eq:two-cluster-bottom} is fixed.
\end{proof}

The extremes agree with the geometry in \cref{sec:conditioning}.  At a
unique nondegenerate optimum---equivalently here, a singleton optimal face
with a nondegenerate optimal basis---$r=p$: if $h\in\ker A$ and $h_N=0$,
then the optimal basis gives $A_Bh_B=0$ and hence $h=0$.  The bottom cluster
is empty and the condition number remains bounded.  If the optimal face has positive
dimension, $\ker(P_NZ)$ contains its tangent space and the bottom cluster is
nonempty; its fixed scale and the $\mu^{-2}$ top scale recover the chord-law
divergence.

The two-window statement is specific to polyhedral log-barrier tails.
The log-det central path of \eqref{eq:degree-two-sdp}, the degree-two SDP
witness in \cref{prop:degree-two-witness}, has four increasingly ordered
reduced-Hessian eigenvalues in the inherited Frobenius metric:
\[
 \lambda_1\sim\sqrt2\,\mu^{-1/2},\qquad
 \lambda_2\sim\mu^{-1},\qquad
 \lambda_3\sim\sqrt2\,\mu^{-3/2},\qquad
 \lambda_4\sim\frac47\mu^{-2}.
\]
To see the metric normalization, use tangent coordinates $(b,g,t,e)$
for the symmetric matrix with diagonal $(1-g-b,g,b)$ and upper
entries $(b,t,e)$. At the center $t=e=0$, write $q=ag-b^2$ and
$A_2=\bigl(\begin{smallmatrix}a&b\\b&g\end{smallmatrix}\bigr)$.
The off-block operator is $A_2^{-1}/b$, giving the first and third
scales because $b\sim\sqrt{\mu/2}$ and $q\sim\mu$.
The $(b,g)$ coordinate Hessian $K$ has Gram matrix
$G=\bigl(\begin{smallmatrix}4&1\\1&2\end{smallmatrix}\bigr)$;
its operator is $G^{-1}K$. Direct differentiation gives
$\mu K_{bb}\to6$, $\mu^{3/2}K_{bg}\to-\sqrt2$, and
$\mu^2K_{gg}\to1$. Thus the larger eigenvalue has coefficient
$(G^{-1})_{22}=4/7$, while
$\mu^3\det K\to4$ and $\det G=7$ give the smaller coefficient one.
These four distinct powers fix the eventual order. In particular, there
is no $O(1)$ window. The complete calculation, including the actual
second derivative, the characteristic polynomial with multiplicities,
and the full central-path tail, is verified in Lean as documented in
\texttt{formal/FRACTIONAL\_SDP.md}. It replaces the earlier numerical
four-window observation with exact asymptotics.

\subsection{Why conjugate gradients see the clusters}
\label{sec:spectra-cg}

Cluster-sensitive CG estimates have a long classical lineage
\cite{axelsson1994-iterative-solution-methods}.  The following elementary
construction is enough here and makes the logarithmic dependence on the
inter-cluster gap explicit.

\begin{proposition}[Two intervals are Krylov-benign]
\label{prop:two-cluster-cg}
Let $H\succ0$ have
\[
 \operatorname{spec}(H)\subseteq[a_1,b_1]\cup[a_2,b_2],
 \qquad 0<a_1\leq b_1<a_2\leq b_2,
\]
and set $\rho_i=b_i/a_i$ and $\Lambda=b_2/a_1$.  Thus
$\kappa(H)\leq\Lambda$, with equality only when both enclosure endpoints
are attained.  For every $0<\epsilon\leq1/2$ there is a residual
polynomial $q$, with $q(0)=1$ and
$\max_{\lambda\in\operatorname{spec}(H)}|q(\lambda)|\leq\epsilon$, of
degree
\begin{equation}
 O\!\left(
  \sqrt{\rho_1\rho_2}\log\Lambda\log(1/\epsilon)
  +\sqrt{\rho_2}\log(1/\epsilon)
 \right).
 \label{eq:two-cluster-cg-degree}
\end{equation}
Consequently CG reaches relative $H$-norm error $\epsilon$ in
$O_\rho(\log(\Lambda/\epsilon)\log(1/\epsilon))$ iterations.
\end{proposition}

\begin{proof}
If $a_2<b_2$, let $p_2$ be the shifted-Chebyshev residual for $[a_2,b_2]$,
normalized by $p_2(0)=1$, of degree
$m_2=\lceil\sqrt{\rho_2}\rceil$.  The standard Chebyshev estimate gives
\[
 \max_{[a_2,b_2]}|p_2|
 \leq2e^{-2m_2/\sqrt{\rho_2}}\leq\frac12.
\]
On $[0,a_2]$ its magnitude is at most one: after the affine map from
$[a_2,b_2]$ to $[-1,1]$, the normalized Chebyshev quotient is monotone in
magnitude between the image of zero and the left endpoint.  At the edge
$a_2=b_2$, instead take $p_2(x)=1-x/a_2$ and $m_2=1$; the same two bounds
hold.

If $a_1=b_1$, take the exact factor $r_1(x)=1-x/a_1$ and $m_1=1$.
Otherwise let $r_1$ be the shifted-Chebyshev residual for $[a_1,b_1]$ of
degree $m_1=O(\sqrt{\rho_1}\log(1/\epsilon))$, with $r_1(0)=1$ and
$\max_{[a_1,b_1]}|r_1|\leq\epsilon$.  In its explicit quotient, the
denominator is $T_{m_1}(\eta_1)$, where
$\eta_1=(a_1+b_1)/(b_1-a_1)>1$, and
$T_{m_1}(\eta_1)\geq\eta_1^{m_1}$.  Combining this denominator bound with
$|T_m(t)|\leq(2|t|)^m$ for the numerator outside the approximation interval
gives, also for the exact-factor edge case,
\begin{equation}
 \max_{[a_2,b_2]}|r_1|\leq(4\Lambda)^{m_1}.
 \label{eq:lower-chebyshev-growth}
\end{equation}
Choose
\[
 j=\left\lceil m_1\log_2(4\Lambda)+\log_2(1/\epsilon)\right\rceil,
 \qquad q=r_1p_2^j.
\]
Then $q(0)=1$.  On the first interval, $|q|\leq\epsilon$ because
$|p_2|\leq1$; on the second, \eqref{eq:lower-chebyshev-growth} and the
choice of $j$ again give $|q|\leq\epsilon$.  Its degree is
$m_1+jm_2$, which is \eqref{eq:two-cluster-cg-degree}.  CG minimizes the
$H$-norm error over residual polynomials of its degree, so its error is no
larger.
\end{proof}

The numerical check used a 40-variable simplex-slice LP with a
two-coordinate active set and a 37-dimensional optimal face.  From
$\mu=10^{-2}$ to $10^{-8}$, $\kappa$ grew from about $10$ to
$7\times10^{12}$, the bottom-cluster ratio stayed below $1.006$, and the top
ratio approached $4.004$ (from $2.44$ at the shallowest point).  CG solved
random right-hand sides to a recursively updated residual tolerance of
$10^{-8}$ in 3--8 iterations.  At $\mu=10^{-8}$, the reported eight-step
recursive residual masks a postcomputed true relative residual
$\norm{b-Hx}_2/\norm b_2\approx10^{-4}$ and relative $H$-norm error about
$6\times10^{-6}$, a float64 loss-of-conjugacy effect.  A separate 80-bit
check, outside the accompanying scripts, used six iterations;
the benign iteration count therefore survives the precision check.  The
unstructured $\sqrt\kappa$ estimate is about $2.6\times10^6$ iterations.
The accompanying analysis scripts produce the float64 spectrum and
recursive-residual counts.

\subsection{Quantum polynomial boundedness and the access-model gap}
\label{sec:spectra-quantum}

CG evaluates its residual polynomial only on eigenvalues.  QSVT has a
different contract: after normalization, its polynomial must be bounded on
all of $[-1,1]$, including the empty spectral gap and the negative half of
the interval.  This difference rules out transplanting the composite
polynomial in \cref{prop:two-cluster-cg}.  An order-of-magnitude evaluation
of this construction on representative intervals reaches about $10^{85}$
inside the gap; this illustrative estimate is not used below.

Normalize $\norm{H}_2=1$ and write
\begin{equation}
 \operatorname{spec}(H)\subseteq
 [1/\Lambda,\rho_1/\Lambda]\cup[1/\rho_2,1],
 \qquad \Lambda=\kappa(H),
 \label{eq:quantum-two-cluster-normalization}
\end{equation}
where $\rho_1,\rho_2\geq1$ are fixed.  Cost statements below
assume norm-tight encodings,
$\alpha_H=\Theta(\norm H)$ and $\alpha_G=\Theta(\norm G)$.  Without this
hypothesis the plain-access parameter is
$\kappa_{\rm BE}=\alpha_H\norm{H^{-1}}$ from \cref{eq:kappa-be}, and the
factor parameter is $\kappa_G=\alpha_G\norm{H^{-1}}^{1/2}$.

\begin{theorem}[Theorem Q: clustered solves in the two access models]
\label{thm:quantum-cluster-obstruction}
For the two-cluster spectra \eqref{eq:quantum-two-cluster-normalization}:
\begin{enumerate}[label=\textup{(\roman*)}]
 \item In the one-sided polynomial model, a residual polynomial bounded by
 $M$ on $[0,1]$, with $q(0)=1$ and
 $|q(1/\Lambda)|\leq\epsilon\leq1/2$, has degree
 \begin{equation}
  d\geq\sqrt{\frac{(1-\epsilon)\Lambda}{2M}}.
  \label{eq:markov-one-sided-no-go}
 \end{equation}
 Thus its degree--amplitude-normalization metric satisfies
 $dM=\Omega(\sqrt\Lambda)$.
 \item In the fixed-polynomial plain-access model, any single polynomial
 that, for $\rho_1\geq2$, relatively inverts a bottom interval
 containing the points $y=1,2$ to accuracy at most $1/8$, and is bounded
 by $M$ on the QSVT-required symmetric interval
 $[-\Lambda,\Lambda]$, obeys
 \begin{equation}
  dM\geq\frac5{16}\sqrt{\Lambda^2-4}=\Omega(\Lambda).
  \label{eq:symmetric-bernstein-no-go}
 \end{equation}
 This is a lower bound on the polynomial metric for a right-hand side
 supported on the bottom cluster, not a general degree-to-QLS-cost theorem.
 \item With factor access $H=G^\dagger G$, the one-sided construction is an
 even singular-value polynomial with
 $dM=\Otilde(\sqrt\Lambda)$.  Item~\textup{(i)}, applied to the residual of
 this construction, gives the matching
 $\Omega(\sqrt\Lambda)$ degree lower bound.
 \item Let $N_{\rm QLS}$ denote the Hilbert-space dimension, distinct from
 the optimal-partition index set $N$ above.  In the full query model,
 plain access has $\Thetatilde(\kappa)$ worst-case QLS cost for
 $N_{\rm QLS}\gtrsim\kappa$, even for two-point clusters.  More precisely,
 \citet[Proposition~6]{orsucci2021-on-solving-classes-of-positive} prove an
 $\Omega(\min\{\kappa,N_{\rm QLS}\})$ normalized-block-encoding lower bound
 for all algorithms.  Their Proposition~7 supplies the sparse/Grover lower
 bound $\Omega(\min\{\kappa,\sqrt{N_{\rm QLS}}\})$, whose
 linear-in-$\kappa$ regime requires $N_{\rm QLS}\geq\kappa^2$.
\end{enumerate}
\end{theorem}

\begin{proof}
For (i), the mean-value theorem gives a point
$\xi\in(0,1/\Lambda)$ with
$|q'(\xi)|\geq(1-\epsilon)\Lambda$.  Markov's inequality, rescaled from
$[-1,1]$ to $[0,1]$, gives
$\max_{[0,1]}|q'|\leq2d^2M$.  Rearranging proves
\eqref{eq:markov-one-sided-no-go}.  Since $M\geq|q(0)|=1$, multiplying the
bound by $M$ gives $dM=\Omega(\sqrt\Lambda)$.

For (ii), use the rescaled eigenvalue $y=\Lambda\lambda$.  Relative error
$\delta'\leq1/8$ for the target $1/y$ implies
\[
 p(1)\geq1-\delta',
 \qquad
 p(2)\leq\frac{1+\delta'}2,
 \qquad
 p(1)-p(2)\geq\frac5{16}.
\]
The mean-value theorem therefore gives a $\xi\in(1,2)$ with
$|p'(\xi)|\geq5/16$.  Bernstein's inequality on the symmetric interval
states
\[
 |p'(y)|\leq\frac{dM}{\sqrt{\Lambda^2-y^2}}.
\]
Since $|\xi|<2$, this yields
\eqref{eq:symmetric-bernstein-no-go}.

The metric $dM$ has the stated operational meaning only after fixing the
input support.  For a right-hand side in the bottom cluster, normalizing a
polynomial bounded by $M$ reduces the useful transformed amplitude by
$\Theta(M)$, so amplitude amplification gives the $dM$ proxy.  For a general
right-hand side the success amplitude also depends on its spectral overlap.
Indeed, on a two-point spectrum the degree-one separator
$p(y)=1-(y-1)/\Lambda$ has $dM=O(1)$ even though the all-algorithms result in
item~\textup{(iv)} still gives an $\Omega(\kappa)$ worst case.  Thus
item~\textup{(ii)} is not the basis for the plain-access row of the cost
table; item~\textup{(iv)} is.

For (iii), the matching lower bound is item~\textup{(i)} applied to the
residual of the construction below.  Writing $q(y)=1-y\,p_y(y)$ in the
rescaled variable $y=\Lambda\lambda$, one has $q(0)=1$,
$|q(1)|\le e^{-T}+O(\epsilon)=O(\epsilon)$, and
$|q(y)|\le e^{-Ty}+y\,|p_y(y)-f_T(y)|\le1+O(\epsilon)$ on
$[0,\max\{L_G,\Lambda\}]\supseteq[0,\Lambda]$, the approximation interval of
the construction below.  In $\lambda$ units this is a residual
polynomial of degree $d+1$ bounded by $1+O(\epsilon)$ on $[0,1]$, so
item~\textup{(i)} forces $d=\Omega(\sqrt\Lambda)$; since
$M\ge|p_y(1)|\ge1-O(\epsilon)$, also $dM=\Omega(\sqrt\Lambda)$.  Here $M$ is
the amplitude bound of the inverse polynomial in $y$ units, whose target is
$1/y$; an inverse approximant in $\lambda$ units necessarily has bound at
least $(1-\epsilon)\Lambda$, so the two normalizations are not
interchangeable.

Now include the encoding normalization explicitly.  A block encoding of
$G$ presents singular values $s=\sigma/\alpha_G$ in $[-1,1]$.  Put
$L_G=\Lambda\alpha_G^2$.  If $p_y$ is bounded on $[0,L_G]$, then
\begin{equation}
 q_G(s)=p_y(\Lambda\alpha_G^2s^2)
 \label{eq:even-factor-polynomial}
\end{equation}
is even and bounded on $[-1,1]$.  On a singular value of $G$, its argument
is $\Lambda\sigma^2$, exactly the one-sided eigenvalue variable for
$G^\dagger G$.  Norm-tightness gives $\alpha_G=\Theta(\norm G)=\Theta(1)$
under \eqref{eq:quantum-two-cluster-normalization}, so
$L_G=\Theta(\Lambda)$; it does not set $\alpha_G$ equal to one.

For completeness, a $\Otilde(\sqrt\Lambda)$ one-sided polynomial follows from
the endpoint geometry.  In $y=\Lambda\lambda$ units define
\[
 f_T(y)=\frac{1-e^{-Ty}}y=\int_0^T e^{-ty}\,dt,
 \qquad T=\Theta(\log(1/\epsilon)).
\]
For $y\geq1$, its relative error as an approximation to $1/y$ is at most
$e^{-T}$, and $|f_T(y)|\leq T$ on $[0,L_G']$, where
$L_G'=\max\{L_G,\Lambda\}=\Theta(\Lambda)$.  The shifted-Chebyshev expansion
of $e^{-ty}$ on this interval, truncated uniformly for
$0\leq t\leq T$, has degree
\[
 O\!\left(\sqrt{L_G'T\log(L_G'T/\epsilon)}
       +\log(L_G'T/\epsilon)\right).
\]
Integrating its coefficients over $t$ produces a polynomial $p_y(y)$ approximating
$f_T$ to absolute error $O(\epsilon/\Lambda)$, hence to relative error
$O(\epsilon)$ for every $y\in[1,\Lambda]$.  Its bound is
$M=O(T)$, so degree times subnormalization is
$\Otilde(\sqrt\Lambda)$ under norm-tight access.  The even polynomial
$q_G$ in \eqref{eq:even-factor-polynomial} has only twice the degree.  This
makes the construction admissible in the factor-access
polynomial model and proves item~\textup{(iii)}.  Algebraically, the
corresponding factor route to the linear-system solution is also
transparent.  For a full-column-rank factor, first form
$w=(G^\dagger)^+b$ and then $x=G^+w$; an SVD verifies
$x=(G^\dagger G)^{-1}b$.  Define the two relative output amplitudes
\[
 \zeta_1=
 \frac{\norm{(G^\dagger)^+b}}{\norm{(G^\dagger)^+}\norm b},
 \qquad
 \zeta_2=
 \frac{\norm{G^+w}}{\norm{G^+}\norm w}.
\]
Each pseudoinverse application costs $\Otilde(\kappa_G)$ to constant
relative amplitude.  The end-to-end
$\Otilde(\kappa_G)=\Otilde(\sqrt\kappa)$ solve follows only under the
explicit overlap hypothesis $\zeta_1,\zeta_2=\Omega(1)$; otherwise amplitude
amplification pays the corresponding inverse overlaps.

There is no generic factor-access QLS guarantee at
$\Otilde(\sqrt\kappa)$.  On the Orsucci--Dunjko two-point instance, the
square-root factor is diagonal and its oracle is implementable with
$O(1)$ hidden-input queries.  A generic square-root-cost factor solver would
therefore contradict their $\Omega(\kappa)$ query lower bound.  The factor
statement here is a polynomial-metric result and an overlap-hypothesis solve,
not an unrestricted QLS upper bound.
Item~\textup{(iv)} is exactly the prior-art all-algorithms grounding for the
plain-access worst case.  Together with the standard
$\Otilde(\kappa_{\rm BE})$ block-encoding QLS upper bound, it establishes
the norm-tight $\Thetatilde(\kappa)$ worst-case row below.
\end{proof}

When $\rho_1=1$, item~\textup{(ii)} is vacuous because its two interpolation
points are absent; this does not make the case algorithmically easy.
Orsucci--Dunjko's all-algorithms lower bound in item~\textup{(iv)} already
covers exact two-point clusters.

The symmetric constraint is decisive.  The polynomial in the last proof is
well behaved on $[0,\Lambda]$ but grows on the negative half; it is not a
plain-access QSVT polynomial.  The value $f_T(0)=T\ne0$ also prevents the odd
extension used for the usual $1/x$ target.  A discretized Remez linear
program makes the distinction visible.  Requiring a polynomial to separate
$p(1/\Lambda)\leq0.1$ from $p(2/\Lambda)\geq0.9$ while staying bounded by one
gave the following minimum degrees; the computation is included in the
accompanying analysis scripts.
\begin{center}
\begin{tabular}{@{}rrrrr@{}}
\toprule
$\Lambda$ & $d_{[-1,1]}$ & $d_{[-1,1]}/\Lambda$
 & $d_{[0,1]}$ & $d_{[0,1]}/\sqrt\Lambda$ \\
\midrule
50  & 52  & 1.04 & 12 & 1.70 \\
100 & 102 & 1.02 & 17 & 1.70 \\
200 & 205 & 1.02 & 24 & 1.70 \\
400 & 376 & 0.94 & 33 & 1.65 \\
\bottomrule
\end{tabular}
\end{center}
Thus the symmetric calculation is $0.94$--$1.04$ times $\Lambda$, whereas
the one-sided relaxation is about $1.7\sqrt\Lambda$.

Suppressing logarithmic precision factors and keeping the access model
fixed, the result is the following solve-cost comparison.
\begin{center}
\begin{tabular}{@{}lll@{}}
\toprule
solver & matrix access & two-cluster cost \\
\midrule
classical CG & matrix--vector products & $\polylog(\kappa)$ \\
factor polynomial route\footnotemark & norm-tight $H=G^\dagger G$
 & $\Thetatilde(\sqrt\kappa)$ \\
QLS worst case, $N_{\rm QLS}\gtrsim\kappa$
 & norm-tight plain encoding of $H$ & $\Thetatilde(\kappa)$ \\
\bottomrule
\end{tabular}
\end{center}
\footnotetext{The factor row is a polynomial-metric statement and an
end-to-end solve cost only under the overlap hypothesis
$\zeta_1,\zeta_2=\Omega(1)$; it is not a generic factor-access QLS
guarantee.}
With \cref{sec:newton-formulations}'s convention
$D=X^{1/2}S^{-1/2}$, the normal-equations matrix is
$H=AD^2A^\top$, and the factor $G=DA^\top$ satisfies
$H=G^\dagger G$ and is naturally available from
the IPM data when the corresponding factor oracle can actually be
implemented.  One should act on that factor rather than
compose it into an encoding of $H$ and inadvertently square the singular
condition number.

The all-algorithms plain-access construction is prior art due to Orsucci and
Dunjko.  Somma and de Wolf's sum-of-squares separation is also prior art,
but it is a bound for guided ground-state energy estimation, under their
dimension condition
$\dim\geq\log(1/\epsilon)/\gamma_{\rm guide}^2$, where
$\gamma_{\rm guide}$ is their guide overlap, and their stated error assumptions
\cite{somma2026-optimal-lower-bound-ground-state-energy}.  It distinguishes
$\Omega(1/\delta)$ plain resolution from
$\Omega(1/\sqrt\delta)$ sum-of-squares resolution; we do not claim a transfer
of that result to QLS.  Within the polynomial model, the matching lower
bound for the even/one-sided factor class is instead item~\textup{(i)}.
Cluster-sensitive Krylov convergence is also classical.  What appears
to be new here is the observation that the hard two-cluster quantum instances
are solved by CG in $O(1)$ iterations when the clusters are points, together
with the resulting separation on central-path Newton spectra.  The novelty
claim is this classical-side observation and assembly, not the quantum lower
bounds themselves.

\begin{openproblem}[Exact-subnormalization shift]
\label{open:exact-subnormalization-shift}
The sharpest access-model question left by Theorem~Q is the following.
The shift $H\mapsto I-H$ moves the bottom cluster to an endpoint, where
$1/\Lambda$-scale discrimination would have square-root polynomial degree.
Can a subnormalization-one block encoding of $I-H$ be constructed from plain
access to $H$ without paying the apparent $\Theta(\Lambda)$ uniform
amplification cost near the new spectral edge?
\end{openproblem}

Orsucci and Dunjko raise the closely related normalization problem for
$B=I-\eta A$ (in their notation) in their Section~4.3
\cite{orsucci2021-on-solving-classes-of-positive}.  They observe that a
linear combination of $I$ and a normalized block encoding of $A$ yields only
the subnormalized $pB$, that any black-box amplification of $B/\beta$ to $B$
is in general inefficient by the lower bound of
\cite[Theorem~73]{andras2019-quantum-singular-value-transformation-beyond}
applied to $f(x)=\beta x$, and they leave open whether a normalized block
encoding of a matrix with operator norm at most one can be built from a
sparse-matrix oracle for it, answering it positively for diagonally dominant
$A$ and for sums of positive semidefinite local Hamiltonians.
\Cref{open:exact-subnormalization-shift} is the promise-sensitive form of
their black-box obstruction: the subnormalization to remove is the fixed
factor two, the spectrum of $H$ is promised to satisfy
\eqref{eq:quantum-two-cluster-normalization}, and the question is whether that
promise lowers the $\Theta(\Lambda)$ uniform-amplification degree at the
shifted edge, not whether a general normalization procedure exists.

\subsection{The central path does not use a worst-case right-hand side}
\label{sec:spectra-rhs}

The preceding obstruction concerns arbitrary right-hand sides.  The exact
log-barrier path generates a much more special vector.  Let
$P_{\rm bot}(\mu)$ denote the spectral projector of $H_\mu$ onto the bottom
cluster in \cref{prop:two-cluster-hessian}.

\begin{proposition}[Universal Newton right-hand side and weak-mode alignment]
\label{prop:newton-rhs-alignment}
In the setting of \cref{prop:two-cluster-hessian}, at an exact central point
define $v_\mu=Z^\top X^{-1}e$.  For a centering target
$\mu_2=\gamma\mu$ with $0<\gamma<1$, the exact reduced Newton equations are
\begin{equation}
 \mu_2H_\mu d_c=r_c,
 \qquad r_c=-(1-\gamma)\mu v_\mu,
 \qquad
 \mu H_\mu d_a=r_a,
 \qquad r_a=-\mu v_\mu.
 \label{eq:universal-newton-rhs}
\end{equation}
Here $d_a$ is the affine-scaling direction formed with the current local
barrier metric; $\gamma$ avoids collision with the singular value $\sigma$
in \cref{prop:two-cluster-hessian}.  If the right-hand side is nonzero,
then
\begin{equation}
 \frac{\norm{P_{\rm bot}(\mu)r}_2}{\norm r_2}=O(\mu^2)
 \qquad(r=r_c\text{ or }r_a).
 \label{eq:rhs-bottom-fraction}
\end{equation}
For either equation in \eqref{eq:universal-newton-rhs}, if the bottom
component of its exact solution is dropped, the resulting relative
Euclidean residual is exactly the fraction in \eqref{eq:rhs-bottom-fraction}.
\end{proposition}

\begin{proof}
At exact centrality,
$Z^\top(c-\mu X^{-1}e)=0$.  Negating the reduced gradient for the target
$\mu_2$ gives
\[
 -Z^\top(c-\mu_2X^{-1}e)
 =-(1-\gamma)\mu Z^\top X^{-1}e,
\]
and setting the target barrier term to zero gives
$-Z^\top c=-\mu Z^\top X^{-1}e$.  This proves
the right-hand sides in \eqref{eq:universal-newton-rhs}.  The reduced Hessian
of the target barrier objective is $\mu_2H_\mu$; the affine direction uses
the current local metric $\mu H_\mu$.  Thus the displayed system scalars are
also exact.

For the alignment claim, recall the splitting
$H_\mu=H_{\mu,N}+H_{\mu,B}$ in \eqref{eq:two-cluster-split}.  The bottom
eigenspace of $H_{\mu,N}$ is exactly
\begin{equation}
 \mathcal T_\Phi=\ker(P_NZ),
 \label{eq:face-tangent-space}
\end{equation}
the reduced tangent space of the primal optimal face.  The positive spectrum
of $H_{\mu,N}$ begins at $\Theta(\mu^{-2})$, whereas
$\norm{H_{\mu,B}}_2=O(1)$.  The Davis--Kahan theorem therefore gives
\begin{equation}
 \norm{P_{\rm bot}(\mu)-P_{\mathcal T_\Phi}}_2=O(\mu^2)
 \label{eq:bottom-projector-rotation}
\end{equation}
\cite{davis1970-rotation-eigenvectors-perturbation}.

Let $x^a$ be the analytic center of the primal optimal face.  The log-barrier
path converges to $x^a$, and strict complementarity gives the tail expansion
$x_B(\mu)=x_B^a+O(\mu)$
\cite{mclinden1980-analogue-moreau-proximation,wright1994-properties-hessian-logarithmic-barrier}.
For every reduced unit vector $u\in\mathcal T_\Phi$, the ambient direction
$Zu$ has $(Zu)_N=0$, and first-order optimality of $x^a$ on the face says
\[
 \inner{(Zu)_B}{(X_B^a)^{-1}e_B}=0.
\]
Since the coordinates of $x_B^a$ are positive, the tail expansion then gives
$\inner{u}{Z^\top X^{-1}e}=O(\mu)$.  On the other hand
$\norm{X^{-1}e}_2=O(1/\mu)$ because $x_N=\mu S_N^{-1}e_N$.
Combining these facts with \eqref{eq:bottom-projector-rotation} yields
\begin{equation}
 \norm{P_{\rm bot}(\mu)v_\mu}_2=O(\mu).
 \label{eq:bottom-rhs-amplitude}
\end{equation}
Centrality also gives $v_\mu=Z^\top c/\mu$.  The nonconstant-objective
hypothesis makes $Z^\top c\ne0$, so $\norm{v_\mu}_2=\Theta(1/\mu)$; dividing
\eqref{eq:bottom-rhs-amplitude} proves \eqref{eq:rhs-bottom-fraction}.

Finally, write $d=d_{\rm top}+d_{\rm bot}$ for either system
$\tau H_\mu d=r$, where $\tau=\mu_2$ for centering and $\tau=\mu$ for the
affine direction, and use $\widetilde d=d_{\rm top}$.  Then
\[
 \tau H_\mu\widetilde d-r
 =-\tau H_\mu d_{\rm bot}=-P_{\rm bot}(\mu)r.
\]
Taking Euclidean norms proves the exact residual identity.
\end{proof}

\begin{corollary}[Condition-number-free top-window solve]
\label{cor:top-window-solve}
Suppose an inexact Newton step only requires a relative linear-system
residual at most $\eta$, with $\eta$ independent of $\mu$.  On the exact
log-barrier path, once $\mu=O(\sqrt\eta)$, it suffices to invert the top
cluster.  Plain-access QSVT can do so with degree
\begin{equation}
 O\!\left(\rho_2\log(\rho_2/\eta)\right),
 \label{eq:top-window-qsvt-degree}
\end{equation}
independent of $\kappa(H_\mu)$.
\end{corollary}

\begin{proof}
Dropping the bottom solution component contributes $O(\mu^2)$ relative
residual by \cref{prop:newton-rhs-alignment}; choose the tail constant so
that this is at most $\eta/2$.  After scaling the top interval to
$[1/\rho_2,1]$, the standard bounded QSVT approximation to $1/x$ has degree
$O(\rho_2\log(\rho_2/\eta))$ and can make the remaining residual at most
$\eta/2$.  The interval stays a fixed distance from zero, so this polynomial
is admissible with plain access and has no dependence on the remote bottom
cluster.
\end{proof}

The contract in this corollary is the reduced-primal relative residual
$\norm{\tau H_\mu\widetilde d-r}/\norm r$, with $\tau=\mu_2$ for centering
and $\tau=\mu$ for the affine direction.  Here $r_c$ denotes the reduced
primal right-hand side in \eqref{eq:universal-newton-rhs}, not the KKT
complementarity residual denoted by $r_c$ in \cref{sec:newton-formulations}.
This contract is distinct from the absolute OSS and MNES conventions in
\cref{def:inexact-newton-residual}: OSS requires
an $O(\bar\mu)$ residual, while MNES requires an
$O(\sqrt{\bar\mu/n})$ system residual chosen to produce an
$O(\bar\mu)$ complementarity residual.  After identifying residual vectors,
a relative tolerance $\eta_{\rm rel}$ implies an absolute requirement
$\norm{\mathrm{res}}\leq\tau_{\rm abs}$ only when
$\eta_{\rm rel}\norm r\leq\tau_{\rm abs}$.  Here $r_a=-Z^\top c$ and
$r_c=-(1-\gamma)Z^\top c$ are generally constant in norm, so a constant
relative tolerance does not by itself imply either shrinking absolute
contract.

The measured law is substantially smaller than its asymptotic bound.  On
Netlib \texttt{afiro}, the dropped-bottom residual was
$6.7\times10^{-6}$, $6.4\times10^{-10}$, and $6.7\times10^{-14}$ at
$\mu=10^{-1},10^{-3},10^{-5}$, respectively: an instance constant of about
$7\times10^{-4}$ multiplying $\mu^2$.  Both centering and affine columns
agree, as \eqref{eq:universal-newton-rhs} predicts; the top-cluster RHS mass
rounded to $1.0000$, and the solution's top mass remained at least $0.9929$.
The symmetric 40-variable family had machine-zero weak-mode mass.  The
accompanying analysis scripts generate these measurements.

An end-to-end damped-Newton run tested more than exact central points.  It
used shrink factor $0.85$, up to four corrections per parameter value, and
fraction-to-boundary steps.  A window solver dropped the bottom modes only
after detecting a gap larger than two decades and otherwise used the full
solve.  The window and exact runs agreed in final gap and correction count:
the toy family finished at gap $1.758\times10^{-8}$ with 205 corrections in
both modes, and \texttt{afiro} finished at relative gap
$7.4\times10^{-10}$ with 393 corrections in both.  This does not establish
iterate-by-iterate identity; direct instrumentation of the two runs shows
that intermediate iterates can differ.  In a separate, independently reproduced experiment, a
naive geometric-mean threshold without gap detection stalled on
\texttt{afiro} at relative gap $0.98$ by splitting the early spectral
continuum arbitrarily.  That negative case is separate from the reported
adaptive experiment.  The operational rule is therefore: use full solves
until a genuine cluster gap opens, then drop the weak modes.  The accompanying
analysis scripts implement the adaptive experiment.

\subsection{A Netlib conditioning survey}
\label{sec:spectra-netlib}

To test how early the asymptotic picture becomes visible, we surveyed 14
small presolved standard-form Netlib LPs.  For each instance, a damped Newton
method followed the log-barrier path in $\ker A$.  We fitted
$\kappa_{\rm red}(g)\asymp g^{-\alpha}$ over the last four grid points,
spanning about three gap decades, counted clusters using gaps larger than
one decade, and ran CG to $10^{-8}$ tolerance with a 4000-iteration cap at
the deepest point.  A second computation solved directional LPs in 30 fixed
random directions over the objective sublevel polytope.  The resulting
$\widehat D(g)$ is a lower estimate of the diameter, not the diameter
itself; fitting $\widehat D(g)\asymp g^\theta$ and comparing with
$\alpha=2-2\theta$ is therefore a consistency check.  It is computationally
independent of the path Hessian eigensolves, but is not an exact second
measurement of the chord-law diameter.

\begin{table}[ht]
\centering
\small
\begin{tabular}{@{}lrrrrrr@{}}
\toprule
instance & $m$ & $n$ & $\alpha$ & $\kappa_{\rm end}$ & clusters & CG its \\
\midrule
adlittle  & 55  & 136 & 2.07 & $1.7\!\times\!10^{13}$ & 2 & 4000 cap \\
afiro     & 9   & 18  & 2.00 & $3.8\!\times\!10^{12}$ & 2 & 20 \\
beaconfd  & 11  & 30  & 0.02 & $1.3\!\times\!10^6$  & 3 & 28 \\
bore3d    & 36  & 65  & 0.03 & $3.4\!\times\!10^{11}$ & 3 & 521 \\
kb2       & 46  & 71  & 0.36 & $9.6\!\times\!10^9$  & 4 & 198 \\
sc105     & 34  & 65  & 0.03 & $9.9\!\times\!10^3$  & 1 & 55 \\
sc205     & 67  & 125 & 0.03 & $7.2\!\times\!10^5$  & 1 & 286 \\
sc50a     & 17  & 33  & 0.01 & $1.5\!\times\!10^3$  & 1 & 20 \\
sc50b     & 15  & 29  & 0.00 & $5.8\!\times\!10^2$  & 1 & 16 \\
scagr7    & 93  & 147 & 1.30 & $8.1\!\times\!10^{12}$ & 2 & 4000 cap \\
scorpion  & 77  & 136 & 0.12 & $1.1\!\times\!10^5$  & 1 & 160 \\
seba      & 10  & 17  & 0.00 & $1.3\!\times\!10^3$  & 2 & 7 \\
brandy    & 107 & 211 & \multicolumn{4}{l}{numerical failure} \\
recipelp  & 63  & 107 & \multicolumn{4}{l}{numerical failure} \\
\bottomrule
\end{tabular}
\caption{Tail conditioning fits and deepest-point solver diagnostics.  The
cluster count uses a one-decade heuristic and should not be read as an exact
spectral partition.}
\label{tab:netlib-conditioning}
\end{table}

Two endpoint classes dominate \cref{tab:netlib-conditioning}.  The
positive-dimensional-face instances \texttt{afiro} and \texttt{adlittle}
have $\alpha\approx2$, while the vertex-like \texttt{sc} instances,
\texttt{seba}, \texttt{scorpion}, \texttt{beaconfd}, and \texttt{bore3d}
have frozen exponents $\alpha\approx0$; their condition-number levels can
still be large.  For
\texttt{afiro}, $\widehat D(g)$ freezes at $126.84$, so $\theta\to0$ predicts
$\alpha=2.000$, matching the path fit $2.00$.  For \texttt{sc50b},
$\widehat D/g$ freezes at $126$, so $\theta\to1$ predicts $\alpha=0$, matching its
condition-number plateau near $582$.

The intermediate \texttt{kb2} slope is a crossover, not evidence for a
third asymptotic class.  Its random-direction estimate of $\theta$ climbs from
$0.08$ to $0.99$ across the available gap decades, the in-the-wild version
of the near-degeneracy plateau in
\cref{cor:near-degeneracy-plateau}; the measured $\alpha=0.36$ lies in the
knee and will tend to zero.  The \texttt{scagr7} local fitted $\theta$ slopes
range from $0.1$ to $0.6$ and may reflect a slower multiscale crossover.  We
claim no fractional asymptotic exponent for a Netlib instance; the proved
fractional example remains the constructed SDP in
\cref{prop:degree-two-witness}.

The instructive counterpoint is \texttt{adlittle}.  At the first float64-safe
probe, its 81 reduced eigenvalues form one cluster with spread ratio
$\kappa=4.3\times10^4$ (about $4.6$ decades).  The middle probe at
$\kappa=1.85\times10^7$ is also one segment, with spread ratio
$1.9\times10^7$ (about $7.3$ decades), and at the deepest safe probe the
spectrum spans nine decades at $\kappa=1.31\times10^9$.  No decade gap
appears at any of these three depths.  CG correspondingly takes 250, 985,
and 2935 iterations;
counts above 81 reflect finite-precision loss of conjugacy.  The deeper
survey row heuristically reports two clusters at
$\kappa=1.7\times10^{13}$, where float64 spectral diagnostics are already
fragile.  This sharpens rather than contradicts
\cref{prop:two-cluster-hessian}: the asymptotic split is real, but the safe
probe establishes only that it has not opened through
$\kappa\approx1.3\times10^9$.  The follow-up calculation is included in
the accompanying analysis scripts.

That wall is itself visible in \texttt{brandy} and \texttt{recipelp}: at
computed condition numbers above roughly $10^{15}$, roundoff produces
negative eigenvalues for a mathematically positive-definite reduced Hessian.
This is consistent with the classical finite-precision analysis of interior
methods \cite{wright1998-ill-conditioning-computational-error-interior}.
Presolve, rank reduction, and Phase-I choices affect the face approached by
the path, so every exponent in \cref{tab:netlib-conditioning} belongs to the
particular presolved representation.  The accompanying analysis scripts
implement the path survey and sublevel fits; the table reports their
deterministic cached-instance run.

\subsection{Scope of the escape}
\label{sec:spectra-scope}

The two-cluster theorem and the $O(\mu^2)$ right-hand-side fraction use the
LP logarithmic barrier, strict complementarity on a sufficiently small tail,
and exact centrality.  For a general barrier, it is plausible that the
barrier's face-center first-order condition produces an analogous alignment,
perhaps only an $O(\mu)$ fraction, but this is a conjecture: neither the
cluster identification nor the required path-limit expansion has been proved
in that generality.  Off-path iterates can acquire weak-mode components.
Nothing here establishes the same rate for predictor--corrector methods,
infeasible-start variants, or unreduced primal--dual systems.  Accordingly,
\cref{cor:top-window-solve} is an exact-path structural statement and the
trajectory experiment is evidence for the stated adaptive rule, not a
general convergence theorem for those variants.
